\clearpage
\section{Synthetic Time-Series Generator}
\label{sec:implemented-generator}

The implementation generates an executable recurrence before generating its
observations. It retains the expression, numerical parameters, lag assignments,
innovation law, and initialization policy, so several trajectories can be drawn
from the same mechanism. This chapter explains how a mechanism becomes a
reproducible dataset: its vocabulary and representation, construction,
execution, and parallel batching. Section~\ref{sec:generator-stability}
then develops the mathematical restrictions used to control the resulting
trajectories. Experimental protocols and numerical comparisons appear in
Sections~\ref{sec:experimental-design} and~\ref{sec:results-discussion}.

\subsection{Motivation for Explicit Autoregressive Mechanisms}
\label{sec:generator-ar-motivation}

Synthetic pretraining motivates a generator that combines diversity with a
retained description of how each series was produced. KernelSynth, introduced
with Chronos, samples sums and products of Gaussian-process kernels and then
draws complete trajectories from the resulting covariance
\citep[Section~4.2]{ansari2024chronos}. Chronos-2 uses it alongside AR, ETS, and
TSI base generators before introducing multivariate dependence
\citep{ansari2025chronos2}. TimePFN's LMC-Synth mixes independent KernelSynth
latent trajectories through a shared linear mixing matrix
\citep{taga2025timepfn}. Toto~2.0 uses the separate \emph{TempoPFN} prior,
which combines several procedural and stochastic families
\citep{moroshan2025tempopfn,khwaja2026toto2}. Section~1.3 describes these
generation procedures and their computational organization.

An explicit autoregressive equation offers a complementary representation:
each dependence is attached to a source series, a lag, and an operation.
For a scalar linear example,
\begin{equation}
 x_t=b+\sum_{\ell=1}^{L}c_{\ell}x_{t-\ell}+\varepsilon_t.
 \label{eq:implemented-ar}
\end{equation}
The nonlinear AST generator replaces the weighted sum with an expression
containing selected lags and interactions. Here $L$ is the largest lag required
by the mechanism; a forecasting model may receive a longer observed context.

The exploratory AR studies motivate this representation by showing strong
one-step approximation for most evaluated KernelSynth and TimePFN latent
series. Section~\ref{sec:ar-coverage-results} presents the measured coverage
and its limitations; the corresponding experimental setup is in
Section~\ref{sec:ar-coverage-protocol}. These studies concern approximation
of observed dynamics, while recovery of the original stochastic mechanism
remains a separate question.

\subsection{Expression Vocabulary and Representation}
\label{sec:generator-vocabulary}

\subsubsection{Operators, Parameters, and Lagged Observations}

The benchmark adopts a fixed vocabulary for writing recurrence mechanisms.
Its terminals are numerical parameters and past observations; addition,
multiplication, and the hyperbolic tangent combine them into expressions.
The symbol $t$ labels the observation being generated. Each update computes
its value from the available lag history and the mechanism's fixed parameters.
Table~\ref{tab:generator-vocabulary} defines the five elements of this language.

\begin{table}[H]
\centering\small
\begin{tabular}{@{}
>{\raggedright\arraybackslash}p{0.20\linewidth}
>{\raggedright\arraybackslash}p{0.25\linewidth}
>{\raggedright\arraybackslash}p{0.46\linewidth}@{}}
\toprule
Element & Representation & Meaning \\
\midrule
Parameter & $P_a$ &
A real coefficient supplied separately, such as $P_0=0.7$. \\
Past observation & $\operatorname{LAG}(j,L_a)$ &
Reads $x_{j,t-\ell}$ after assigning the positive integer $\ell$ to $L_a$. \\
Addition & $\operatorname{ADD}(e_1,e_2)$ &
Combines two expressions as $e_1+e_2$. \\
Multiplication & $\operatorname{MUL}(e_1,e_2)$ &
Forms $e_1e_2$: coefficients, interactions, and repeated factors. \\
Hyperbolic tangent & $\operatorname{TANH}(e)$ &
Nonlinear; output in $(-1,1)$ and derivative magnitude $\leq1$. \\
\bottomrule
\end{tabular}
\caption{Selected expression vocabulary. Parameters and lagged observations
are leaves, addition and multiplication have two children, and $\tanh$ has
one child. The source index $j$ selects the observed series.}
\label{tab:generator-vocabulary}
\end{table}

\subsubsection{Abstract Syntax Trees and Parameter Binding}
An \emph{abstract syntax tree} (AST) describes operations and their
arguments; a binding supplies numerical coefficients and positive integer
lags. Structure and numerical values are retained separately. Repeated
trajectories preserve both the tree and its binding while resampling initial
conditions or innovations according to the selected policy.

Let $\mathcal E$ denote any valid expression, including a complete output
tree or one of its subtrees. The recursive grammar is
\begin{equation}
\begin{aligned}
\mathcal E ::= {}& P_a \mid \operatorname{LAG}(j,L_a)\\
 &\mid\operatorname{ADD}(\mathcal E_1,\mathcal E_2)
  \mid\operatorname{MUL}(\mathcal E_1,\mathcal E_2)\\
 &\mid\operatorname{TANH}(\mathcal E_1).
\end{aligned}
\label{eq:generator-vocabulary-grammar}
\end{equation}
Here $\mathcal E_1$ and $\mathcal E_2$ are expressions built by the same rules.
The tree stores named coefficient slots $P_a$ and lag slots $L_a$; a binding
assigns real coefficients and positive integer delays. For example,
$P_0\operatorname{LAG}(0,L_0)+P_1$ becomes
$0.7x_{0,t-2}+0.01$ when $P_0=0.7$, $L_0=2$, and $P_1=0.01$.
Negative coefficients express subtraction, and repeated lag factors express
powers such as $x_{0,t-1}^2$. One output tree is retained per target series;
its lag leaves can refer to itself or another series in the system.

\paragraph{Temporal evolution through recurrence.}
Growth can arise from repeated updates of the state. For the deterministic
example
\begin{equation}
 x_t=x_{t-1}+7,\qquad t\geq1,
 \label{eq:generator-recurrent-trend}
\end{equation}
each update adds seven to the previous value. Starting from $x_0$, the
sequence is $x_0,x_0+7,x_0+14,\ldots$. Thus a trend can be represented using
only a lag leaf, a constant, and addition. This example has lag derivative
equal to one: the strict contraction conditions in
Section~\ref{sec:generator-stability} select a narrower set of mechanisms.
The shared vocabulary describes expressions; each generation family adds
its own admissibility conditions.

\paragraph{Innovations, initialization, and method-specific subsets.}
The AST represents the deterministic part of the recurrence. The executor
adds the innovation $\varepsilon_{k,t}$ according to the selected noise law,
and the initial history supplies the observations required by the lags.
These choices are recorded alongside the expression. The original sampler
and bounded-domain constructor use sums and products. The Shao--Wu
construction (Section~\ref{sec:implemented-shao-wu}) wraps lag leaves in
$\tanh$ within nonlinear products, or applies $\tanh$ once after each
complete product, according to the selected mode. The first choice checks
global moment contraction; the second checks a uniform amplitude moment
bound. Heaps
(Section~\ref{sec:implemented-heaps}) expresses its linear VAR coefficients
using parameters, lag leaves, sums, and products between coefficients and
observations.

\paragraph{Alignment with the current implementation.}
The corresponding node classes and bindings are defined in
\path{src/tsfmxai/representations/expression.py}. To restrict the existing
base sampler to the selected terminals, set \texttt{time\_probability} to
zero and \texttt{require\_time} to \texttt{False}.
The current code also retains a \texttt{Time} node for earlier experiments;
the vocabulary adopted here selects the five elements in
Table~\ref{tab:generator-vocabulary}. The controlled-family comparison in
Section~\ref{sec:generator-evaluation-protocol} already uses this restriction.

\subsection{Mechanism Construction}
\label{sec:generator-expression-pipeline}

\citet[Section~2.1]{dascoli2022deep} generate recurrent-sequence examples by
sampling an expression tree, selecting initial conditions, and executing the
recurrence. The implementation adopts this expression-first organization and
retains separate representations of structure and numerical values, as
defined in Section~\ref{sec:generator-vocabulary}. A user-provided equation
and a sampled expression can therefore enter the same execution pipeline.

\begin{figure}[H]
\centering
\begin{tikzpicture}[box/.style={draw,rounded corners,align=center,
    text width=3.6cm,minimum height=1.5cm,font=\small}]
\node[box] (tree) at (0,0) {User formula or\\sampled expression};
\node[box] (bound) at (4.7,0) {AST + coefficients\\+ lag bindings};
\node[box] (family) at (9.4,0) {Family rules\\+ certificate};
\node[box] (paths) at (9.4,-2.1) {Initial history\\+ innovations\\$\longrightarrow$ trajectories};
\node[box] (batch) at (4.7,-2.1) {Context/future windows\\+ mechanism IDs};
\draw[->] (tree) -- (bound);
\draw[->] (bound) -- (family);
\draw[->] (family) -- (paths);
\draw[->] (paths) -- (batch);
\end{tikzpicture}
\caption{Implemented generation stages. The selected family determines the
admissible expressions, initialization, and innovation law. The mechanism
record is retained alongside the observations.}
\label{fig:implemented-generator-pipeline}
\end{figure}

\subsubsection{Random Tree Sampling}
Within the selected vocabulary in Section~\ref{sec:generator-vocabulary},
the base sampler draws a binary-operator count, labels leaves as parameters
or lagged series, and recursively constructs the tree. Its implementation
is in \path{src/tsfmxai/generators/random_expression.py}.
For a subtree containing $o$ binary operators, the
probability of assigning $a$ operators to the left child is proportional to
the number of left shapes times the number of right shapes. These counts
are Catalan numbers. This produces a uniform distribution over ordered full
binary shapes conditional on $o$. Internal nodes
are labelled addition or multiplication. Zero-operator trees consist of one
leaf; required leaf types must fit the requested size.

The current random base grammar therefore uses binary operations; the
Shao--Wu transformation in Section~\ref{sec:implemented-shao-wu}
adds unary \texttt{Tanh} nodes to the same AST
representation. The wider unary--binary grammar in D'Ascoli et al.
provides the precedent. Our operator probabilities, parameter ranges, forced
lag presence, and multivariate source indices define the project's sampling
prior. Slot names are assigned in preorder, preserving the sampled left and
right children. This keeps numerical bindings and serialization reproducible
for a fixed seed.

\subsubsection{User-Provided Equations}
A researcher can specify the expression structure and its bindings directly
within the vocabulary in Section~\ref{sec:generator-vocabulary}. For example,
the skeleton $P_0x_0(t-L_0)+P_1$ becomes
$0.7x_0(t-2)+0.01$ after binding its slots. Numerical parameters and lag
assignments are retained with the AST so the resulting formula is executable
and can be inspected independently of its generated observations.
The selected family then applies its construction and admissibility rules
as described in Section~\ref{sec:generator-stability}; the resulting,
possibly transformed equation is the mechanism used for execution and
ground-truth extraction.

\subsection{Trajectory Generation}
\label{sec:generator-trajectories}

\subsubsection{Initial Histories and Random Innovations}
The initial history supplies the observations required by every lag leaf.
The AST represents the deterministic update; the executor adds the
innovation specified by the mechanism's noise law. These choices are retained
with the equation, and each generation family declares its own initialization
and innovation policies in Section~\ref{sec:generator-stability}.
The initial lag values and recursively generated observations have distinct
roles: the history is the starting state, and context/future windows select
their observations from the generated portion of the returned trajectory.

\subsubsection{Synchronous Multivariate Updates}
At target time $t$, a lag leaf reads $x_{j,t-\ell}$ with $\ell\geq1$.
Equivalently, from the common origin $\tau$, the first generated target uses
$t=\tau+1$. The shared executor
\path{src/tsfmxai/generators/stochastic.py} evaluates every target against the
same previous multivariate history, adds the supplied innovation, and then
appends all new values together. This synchronous update preserves the
intended lag relationships between channels. Model inputs and future targets
can subsequently be separated without changing the recorded equation.

\subsubsection{Multiple Trajectories from One Mechanism}
The equation structure, bound coefficients, and lag assignments stay fixed
while repeated trajectories receive separate initialization and innovation
streams. This distinguishes variation within one mechanism from variation
between sampled equations. The mechanism identifier connects those
trajectories to their shared equation and generation policy. Its executed
AST is also the source for the explanation references developed in
Section~\ref{sec:ground-truth-generation}.

\subsection{Reproducible and Parallel Generation}
\label{sec:generator-parallel}

\subsubsection{Seeds and Independent Random Streams}
Seeds are derived from the root seed and explicit job, attempt, and stream
indices. Local random-number generators isolate these draws from the
caller's global random state. Separate streams identify the mechanism,
initialization, and innovations; changing the number of workers preserves
the indexed draws and their final ordering. Worker-count and seed
independence are tested in the test suite.

\subsubsection{Parallel Jobs and Batch Construction}
The existing batch builder assigns one mechanism and several trajectories to
each job, then assembles windows with a requested minimum number of distinct
skeletons. It is implemented in
\path{src/tsfmxai/datasets/parallel_batch.py}. Ordering is restored when jobs
finish. This batch builder currently uses the original sampler with acceptance
checks; the four-family diagnostic comparison in
Section~\ref{sec:generator-evaluation-protocol} evaluates every attempted
mechanism and records failures. The generator policies and the empirical
comparison therefore retain their respective acceptance and accounting rules.

\subsection{Generation Overview}
\label{sec:generation-overview}

A mechanism is the equation executed at every generated step. Its structure,
coefficients, and lag assignments remain fixed while the initial history and
innovation sequence determine a particular trajectory. The same final equation,
initial history, and numerical innovations produce the same trajectory, up to
floating-point effects. A construction method can change that equation before
execution. The original and transformed equations must therefore be retained
separately, with ground-truth explanations referring to the executed mechanism.

Tables~\ref{tab:generation-overview-linear} and~\ref{tab:generation-overview-nonlinear}
show nine sampled examples, selected under restricted sampling configurations
to make these changes readable. The first column gives the original recurrence;
the middle columns give the recurrence produced by each construction; the last
column explains the feature being illustrated. The cases cover individual lags,
an offset, multiple memory scales, repeated factors, and nonlinear interactions.
They are worked examples rather than estimates of how often those structures
occur under the benchmark's full sampling prior.

\paragraph{Reading the equations and the innovations.}
To keep the tables compact, write $z_\ell=x_{t-\ell}$ and
$T_\ell=\tanh(z_\ell)$. A squared lag is displayed as $z_1^2$; its AST contains
two lag leaves joined by multiplication. The index $t$ is the current generated
step, so the first update after history ending at time zero uses $t=1$.
Every recurrence explicitly includes a fresh additive innovation, identified
by its superscript: O denotes the original AST, J the Jacobian method,
F factor saturation, P product-output saturation, and H Heaps. This term is
the generator's random innovation; forecast errors are measured later by
comparing predictions with generated observations.

Here the configured standard deviation is $\sigma_\varepsilon=0.05$.
The O, F, P, and H policies use i.i.d. Gaussian innovations with distribution
$\mathcal N(0,0.05^2)$. The J policy uses i.i.d. uniform innovations on
$[-\sqrt{3}(0.05),\sqrt{3}(0.05)]$ to respect its finite noise bound while
matching that standard deviation. These method labels identify policies;
actual noise sequences may be shared or independently drawn according to the
comparison protocol. Matching standard deviations alone leaves different noise
laws. A trajectory comparison must also specify its initial history and which
numerical innovations are shared. The constructors scale only the deterministic
expression; the innovation is added afterwards at every step.

\paragraph{Construction settings used in the tables.}
The Jacobian calculation uses $r_D=1$, $q_w=1.25$, target gain 0.9, and safety
multiplier 0.95. It scales the complete AST to meet both the sensitivity and
domain-containment checks of Section~\ref{sec:implemented-domain-contraction}.
The two Shao--Wu-inspired modes use target gain 0.8 and safety multiplier 0.95.
Factor saturation replaces lag leaves inside nonlinear products by their
hyperbolic tangents, then checks global sensitivity using the condition of
\citet[Theorem~5.1]{shao2007spectral}. Product-output saturation bounds each
complete nonlinear product and checks the project's second-moment envelope,
as detailed in Section~\ref{sec:implemented-shao-wu}. Their verification records
therefore certify different properties. Both preserve standalone linear terms
before applying their root multiplier. Even a candidate already within a
target bound receives the implementation's 0.95 safety multiplier.

\paragraph{How the Heaps column is matched to a linear example.}
Heaps constructs a linear VAR from free parameter matrices and an innovation
covariance \citep[Sections~2.1--2.2 and Appendix]{heaps2023stationarity}.
For this illustration, each scalar coefficient in the original linear formula
is supplied as the corresponding one-by-one free matrix $\mathbf U_\ell$;
missing lags receive zero. This is an explicit comparison convention, with
the same numerical entries interpreted in different parameter spaces.
The constructor then reconstructs the final AR coefficients jointly. We choose
its stationary mean to preserve the original intercept: after reconstruction,
divide that intercept by one minus the sum of the final AR coefficients.
With one lag, free entry 2 consequently produces coefficient
$2/\sqrt{5}\approx0.894427$. In R4, the free entries at lags 1 and 3 also
produce a coefficient at lag 2. Nonlinear rows indicate the linear scope of
this constructor. Section~\ref{sec:implemented-heaps} explains its standard
stationary-history sampling; the equations below describe the transition rule
separately from that initialization choice.

\paragraph{Reproduction and interpretation.}
The input equations are actual outputs of
\path{src/tsfmxai/generators/random_expression.py}, with the seeds shown in the
tables. Their configurations use one series, no time terminals, and no automatic
lag-coefficient gates; selected parameter magnitudes are fixed at 0.5 or 2 to
keep the arithmetic readable. Operator counts, sampling weights, and lag ranges
for each case are recorded in \path{scripts/check_generation_overview.py}.
Running \texttt{python} \path{scripts/check_generation_overview.py} reconstructs the
nine samples, checks all 27 AST transformations and five matched Heaps
constructions, and verifies the formulas numerically. Coefficients shown with
$\approx$ are rounded to six decimals; calculations retain full precision.
These examples compare construction rules. Explosion rates, initialization
effects, and generation speed require the separate trajectory experiments.

\begin{landscape}
\begingroup
\renewcommand{\baselinestretch}{1}
\footnotesize
\setlength{\tabcolsep}{3pt}
\setlength{\LTcapwidth}{.97\linewidth}
\renewcommand{\arraystretch}{1.2}
\begin{longtable}{@{}
>{\raggedright\arraybackslash}p{.13\linewidth}
>{\raggedright\arraybackslash}p{.14\linewidth}
>{\raggedright\arraybackslash}p{.15\linewidth}
>{\raggedright\arraybackslash}p{.17\linewidth}
>{\raggedright\arraybackslash}p{.17\linewidth}
>{\raggedright\arraybackslash}p{.19\linewidth}@{}}
\caption{Generation Overview: linear and affine sampled recurrences. Each cell
includes the innovation added at step $t$. The Heaps column uses the explicitly
declared free-parameter convention.}\label{tab:generation-overview-linear}\\
\toprule
Sampled equation / original AST & Jacobian: scaled AST & Shao--Wu: factors & Shao--Wu: product output & Heaps: linear construction & Why this example?\\
\midrule
\endfirsthead
\toprule
Sampled equation / original AST & Jacobian: scaled AST & Shao--Wu: factors & Shao--Wu: product output & Heaps: linear construction & Why this example?\\
\midrule
\endhead
\textbf{Operation}\newline
Keep the sampled formula and add the innovation at each step. &
Scale the complete deterministic expression to meet contraction and domain
bounds. In R2, multiply by $0.4275$. &
Apply $\tanh$ to lag factors inside nonlinear products, then scale the whole
expression to limit sensitivity. In R2, multiply by $0.38$. &
Apply $\tanh$ after the complete nonlinear product, then scale the whole
expression to satisfy the second-moment bound. &
Map free parameters to stable AR coefficients. For AR(1),
$a=u/\sqrt{1+u^2}$; $u=2$ gives $a\approx0.894427$. &
Each following row applies these operations to one input equation.
The innovation is added after the deterministic transformation.\\
\midrule
\textbf{R1; seed 0}\newline
$x_t=z_1+\varepsilon_t^{(\mathrm O)}$ &
$x_t=0.855z_1$\newline $+\varepsilon_t^{(\mathrm J)}$ &
$x_t=0.76z_1$\newline $+\varepsilon_t^{(\mathrm F)}$ &
$x_t=0.76z_1$\newline $+\varepsilon_t^{(\mathrm P)}$ &
$x_t\approx0.707107z_1$\newline $+\varepsilon_t^{(\mathrm H)}$ &
Unit feedback: the original is a random walk. All three constructions change
its memory strength.\\
\addlinespace
\textbf{R2; seed 0}\newline
$x_t=2z_1+\varepsilon_t^{(\mathrm O)}$ &
$x_t=0.855z_1$\newline $+\varepsilon_t^{(\mathrm J)}$ &
$x_t=0.76z_1$\newline $+\varepsilon_t^{(\mathrm F)}$ &
$x_t=0.76z_1$\newline $+\varepsilon_t^{(\mathrm P)}$ &
$x_t\approx0.894427z_1$\newline $+\varepsilon_t^{(\mathrm H)}$ &
Amplifying feedback. Scaling maps R1 and R2 to identical J/F/P drifts, while
the Heaps parameter map retains a difference.\\
\addlinespace
\textbf{R3; seed 2}\newline
$x_t=z_1-0.5$\newline $+\varepsilon_t^{(\mathrm O)}$ &
$x_t\approx0.578485$\newline $\times(z_1-0.5)$\newline $+\varepsilon_t^{(\mathrm J)}$ &
$x_t=0.76(z_1-0.5)$\newline $+\varepsilon_t^{(\mathrm F)}$ &
$x_t=0.76(z_1-0.5)$\newline $+\varepsilon_t^{(\mathrm P)}$ &
$x_t\approx0.707107z_1$\newline $-0.5+\varepsilon_t^{(\mathrm H)}$ &
Offset: the AST root multiplier scales the intercept too. J reserves space
for offset and noise; H preserves the intercept by the stated convention.\\
\addlinespace
\textbf{R4; seed 0}\newline
$x_t=z_1+z_3$\newline $+\varepsilon_t^{(\mathrm O)}$ &
$x_t\approx0.333659$\newline $\times(z_1+z_3)$\newline $+\varepsilon_t^{(\mathrm J)}$ &
$x_t=0.38(z_1+z_3)$\newline $+\varepsilon_t^{(\mathrm F)}$ &
$x_t=0.38(z_1+z_3)$\newline $+\varepsilon_t^{(\mathrm P)}$ &
$x_t\approx0.707107z_1$\newline $-0.5z_2+0.707107z_3$\newline $+\varepsilon_t^{(\mathrm H)}$ &
Separated lags: J weights older memory more strongly. Joint H reconstruction
introduces a lag-2 coefficient.\\
\addlinespace
\textbf{R5; seed 0}\newline
$x_t=z_{12}+\varepsilon_t^{(\mathrm O)}$ &
$x_t\approx0.073444z_{12}$\newline $+\varepsilon_t^{(\mathrm J)}$ &
$x_t=0.76z_{12}$\newline $+\varepsilon_t^{(\mathrm F)}$ &
$x_t=0.76z_{12}$\newline $+\varepsilon_t^{(\mathrm P)}$ &
$x_t\approx0.707107z_{12}$\newline $+\varepsilon_t^{(\mathrm H)}$ &
Long delay: the J weight $1.25^{11}$ requires substantially more attenuation
than in R1. Formula size alone misses this difficulty.\\
\bottomrule
\end{longtable}
\endgroup
\end{landscape}

\begin{landscape}
\begingroup
\renewcommand{\baselinestretch}{1}
\footnotesize
\setlength{\tabcolsep}{3pt}
\setlength{\LTcapwidth}{.97\linewidth}
\renewcommand{\arraystretch}{1.2}
\begin{longtable}{@{}
>{\raggedright\arraybackslash}p{.13\linewidth}
>{\raggedright\arraybackslash}p{.14\linewidth}
>{\raggedright\arraybackslash}p{.15\linewidth}
>{\raggedright\arraybackslash}p{.17\linewidth}
>{\raggedright\arraybackslash}p{.17\linewidth}
>{\raggedright\arraybackslash}p{.19\linewidth}@{}}
\caption{Generation Overview: nonlinear sampled recurrences.
Here $z_\ell=x_{t-\ell}$ and $T_\ell=\tanh(z_\ell)$.
Factor mode verifies global sensitivity; product-output mode verifies a
second-moment bound.}\label{tab:generation-overview-nonlinear}\\
\toprule
Sampled equation / original AST & Jacobian: scaled AST & Shao--Wu: factors & Shao--Wu: product output & Heaps: linear construction & Why this example?\\
\midrule
\endfirsthead
\toprule
Sampled equation / original AST & Jacobian: scaled AST & Shao--Wu: factors & Shao--Wu: product output & Heaps: linear construction & Why this example?\\
\midrule
\endhead
\textbf{Operation}\newline
Keep the sampled formula and add the innovation at each step. &
Scale the complete deterministic expression to meet contraction and domain
bounds. In R2, multiply by $0.4275$. &
Apply $\tanh$ to lag factors inside nonlinear products, then scale the whole
expression to limit sensitivity. In R2, multiply by $0.38$. &
Apply $\tanh$ after the complete nonlinear product, then scale the whole
expression to satisfy the second-moment bound. &
Map free parameters to stable AR coefficients. For AR(1),
$a=u/\sqrt{1+u^2}$; $u=2$ gives $a\approx0.894427$. &
Each following row applies these operations to one input equation.
The innovation is added after the deterministic transformation.\\
\midrule
\textbf{R6; seed 0}\newline
$x_t=z_1z_2+\varepsilon_t^{(\mathrm O)}$ &
$x_t=0.38z_1z_2$\newline $+\varepsilon_t^{(\mathrm J)}$ &
$x_t=0.38T_1T_2$\newline $+\varepsilon_t^{(\mathrm F)}$ &
$x_t=0.95\tanh(z_1z_2)$\newline $+\varepsilon_t^{(\mathrm P)}$ &
Outside the linear VAR vocabulary. &
Two distinct lag factors: contrast a scaled polynomial, bounded factors,
and a bounded complete product.\\
\addlinespace
\textbf{R7; seed 0}\newline
$x_t=z_1^2+\varepsilon_t^{(\mathrm O)}$ &
$x_t=0.4275z_1^2$\newline $+\varepsilon_t^{(\mathrm J)}$ &
$x_t=0.38T_1^2$\newline $+\varepsilon_t^{(\mathrm F)}$ &
$x_t=0.95\tanh(z_1^2)$\newline $+\varepsilon_t^{(\mathrm P)}$ &
Outside the linear VAR vocabulary. &
Repeated coordinate: two leaves refer to one lag. The conservative
derivative bound counts both occurrences.\\
\addlinespace
\textbf{R8; seed 5}\newline
$x_t=z_1^2z_2z_3$\newline $+\varepsilon_t^{(\mathrm O)}$ &
$x_t\approx0.177662$\newline $\times z_1^2z_2z_3$\newline $+\varepsilon_t^{(\mathrm J)}$ &
$x_t=0.19T_1^2T_2T_3$\newline $+\varepsilon_t^{(\mathrm F)}$ &
$x_t=0.95$\newline $\times\tanh(z_1^2z_2z_3)$\newline $+\varepsilon_t^{(\mathrm P)}$ &
Outside the linear VAR vocabulary. &
Four factors: increasing product degree raises the sensitivity budget.
P saturates the entire connected product once.\\
\addlinespace
\textbf{R9; seed 1438}\newline
$x_t=0.5z_1+z_1z_2$\newline $+\varepsilon_t^{(\mathrm O)}$ &
$x_t\approx0.310909$\newline $\times(0.5z_1+z_1z_2)$\newline $+\varepsilon_t^{(\mathrm J)}$ &
$x_t=0.304$\newline $\times(0.5z_1+T_1T_2)$\newline $+\varepsilon_t^{(\mathrm F)}$ &
$x_t=0.95$\newline $\times[0.5z_1+\tanh(z_1z_2)]$\newline $+\varepsilon_t^{(\mathrm P)}$ &
Outside the linear VAR vocabulary. &
Mixed linear and nonlinear feedback: selective saturation preserves the
linear lag; the root multiplier still scales both terms.\\
\bottomrule
\end{longtable}
\endgroup
\end{landscape}


% Retained writing notes for later expansion of this chapter:
% - Document accepted formula syntax, parser validation, and user-facing help.
% - Add a hand-computed indexing example and the context/future array layout.
% - Illustrate text, AST, infix formula, and serialization together.
% - Detail the frontend and mechanism-by-column/trajectory-by-row plots when
%   stable report figures are selected.
% - Expand batching/API coverage when the controlled families are integrated;
%   document DataLoader integration once implemented.
% - Keep mathematical admissibility in Chapter~\ref{sec:generator-stability}
%   and measured generator comparisons in Chapter~\ref{sec:results-discussion}.
